\chapter{Limitaciones conocidas}\label{cap:limitaciones}

\section{Velocidad}

El decodificador va entre 0.5 y 0.8 veces tiempo real según la resolución.
El límite está en la etapa de compensación de movimiento del núcleo, no en la
memoria ni en el reloj (medido en el proyecto). Sirve para clips grabados; no
para fuentes en vivo que no acepten control de flujo.

\section{Matrices de cuantización cargadas}

Los streams que cargan en el encabezado de secuencia una matriz de
cuantización intra no simétrica se decodifican con errores visibles desde la
primera imagen. Afecta a 11 de los 39 streams de conformidad del banco de
pruebas, entre ellos varios que cargan la matriz \emph{por defecto}. Las
matrices simétricas, y el caso habitual en que el stream no carga ninguna,
funcionan correctamente. Es una diferencia del núcleo original, en
investigación.

\section{Deriva en imágenes predichas}

En imágenes P y B con mucho movimiento la diferencia con el decodificador de
referencia crece levemente a lo largo del GOP (error medio absoluto de hasta
$\sim$1 en los streams típicos), sin efecto visible apreciable. También es
propia del núcleo original.

\section{Otras}

\begin{itemize}
  \item Sin salida de video: la placa no tiene conector de video; la salida es
    la red.
  \item MPEG-1 no está soportado; un stream \emph{stuffing} de conformidad
    (\cmd{tcela-15}) no produce imágenes.
  \item Con TLS la entrega queda limitada a unos 15 cuadros por segundo a
    704×480.
  \item Cambios de resolución dentro de un mismo clip: soportados por el
    protocolo, sin verificación exhaustiva en hardware.
\end{itemize}
